% Part: lambda-calculus
% Chapter: lambda-definability
% Section: truth-values

\documentclass[../../../include/open-logic-section]{subfiles}

\begin{document}

\olfileid[id]{lam}{ldf}{tvr}
\olsection{Nilai Kebenaran dan Relasi}

Kita dapat mengodekan nilai kebenaran dalam kalkulus lambda murni sebagai
berikut:
\begin{align*}
  \fn{true} & \ident \lambd[x][\lambd[y][x]]\\
  \fn{false} & \ident \lambd[x][\lambd[y][y]]
\end{align*}

Nilai kebenaran direpresentasikan sebagai \emph{penyeleksi}, yakni fungsi yang
menerima dua argumen dan mengembalikan salah satunya. Nilai kebenaran
$\fn{true}$ memilih argumen pertamanya, sedangkan $\fn{false}$ memilih argumen
keduanya. Sebagai contoh, $\fn{true}\, M N$ selalu tereduksi menjadi $M$,
sementara $\fn{false}\, M N$ selalu tereduksi menjadi~$N$.

\begin{defn}
Kita menyebut suatu relasi $R \subseteq \Nat^n$ !!{lambda definable} jika
terdapat suatu suku~$R$ sedemikian sehingga
\begin{align*}
  R\, \num{n_1} \dots \num{n_k} & \bred \fn{true}
  \intertext{setiap kali $R(n_1, \dots, n_k)$ berlaku, dan}
  R\, \num{n_1} \dots \num{n_k} & \bred \fn{false}
\end{align*}
dalam kasus lainnya.
\end{defn}

Sebagai contoh, relasi $\fn{IsZero} = \{0\}$, yang berlaku pada $0$ dan hanya
pada $0$, bersifat !!{lambda definable} oleh
\[
  \fn{IsZero} \ident \lambd[n][n (\lambd[x][\fn{false}])\, \fn{true}].
\]
Bagaimana cara kerjanya? Karena numeral Church didefinisikan sebagai iterator
(fungsi yang menerapkan argumen pertamanya pada argumen kedua sebanyak $n$
kali), kita menetapkan nilai awal sebagai $\fn{true}$ dan, pada setiap langkah
iterasi, mengembalikan $\fn{false}$ tanpa memedulikan hasil iterasi sebelumnya.
Langkah ini akan diterapkan pada nilai awal sebanyak $n$ kali, dan hasilnya
adalah $\fn{true}$ jika dan hanya jika langkah tersebut sama sekali tidak
diterapkan, yakni ketika $n = 0$.

Berdasarkan representasi nilai kebenaran ini, selanjutnya kita dapat
mendefinisikan beberapa fungsi kebenaran. Berikut dua di antaranya, yaitu
representasi negasi dan konjungsi:
\begin{align*}
  \fn{Not} & \ident \lambd[x][x\, \fn{false}\, \fn{true}]\\
  \fn{And} & \ident \lambd[x][\lambd[y][xy \,\fn{false}]]
\end{align*}
Fungsi ``$\fn{Not}$'' menerima satu argumen dan mengembalikan $\fn{true}$ jika
argumennya adalah $\fn{false}$, serta $\fn{false}$ jika argumennya
adalah~$\fn{true}$. Fungsi ``$\fn{And}$'' menerima dua nilai kebenaran sebagai
argumen dan seharusnya mengembalikan $\fn{true}$ jika dan hanya jika kedua
argumennya adalah~$\fn{true}$. Nilai kebenaran direpresentasikan sebagai
penyeleksi (seperti dijelaskan di atas), sehingga ketika $x$ adalah nilai
kebenaran dan diterapkan pada dua argumen, hasilnya adalah argumen pertama jika
$x$ adalah $\fn{true}$ dan argumen kedua dalam kasus lainnya. Sekarang
$\fn{And}$ menerima dua argumen $x$ dan $y$, lalu memberikan $y$ dan
$\fn{false}$ kepada argumen pertamanya, yaitu~$x$. Dengan mengandaikan bahwa
$x$ adalah nilai kebenaran, hasilnya akan terevaluasi menjadi~$y$ jika $x$
adalah $\fn{true}$, dan menjadi $\fn{false}$ jika $x$ adalah $\fn{false}$,
tepat seperti yang diinginkan.

Perhatikan bahwa di sini kita mengandaikan hanya nilai kebenaran yang digunakan
sebagai argumen bagi $\fn{And}$. Jika fungsi itu diberi suku lain, hasilnya
(yakni bentuk normalnya, jika ada) mungkin bukan nilai kebenaran.

\begin{prob}
  Definisikan fungsi $\fn{Or}$ dan $\fn{Xor}$ yang merepresentasikan fungsi
  kebenaran disjungsi inklusif dan eksklusif dengan menggunakan pengodean nilai
  kebenaran sebagai suku-$\lambd$.
\end{prob}

\end{document}
